Вследствие энергетического перехода во всём мире увеличивается интерес к малым формам гидроэнергетики (МГЭС). 
Повышение энергоэффективности и расширение использования возобновляемых источников энергии (ВИЭ) определены стратегией стран
как основные направления при создании инновационного энергетического сектора. Одним из таких решений является сегнерово колесо.
Оно представляет собой двигатель, основанный на реактивном действии вытекающей воды, первая в истории гидравлическая турбина. 
Его принцип основан на использовании реакции струи жидкости при вытекании из изогнутых трубок. Оно было предложено 
Я.~А.~Сегнером ещё в середине~XVIII века и считалось прибором для демонстрации законов физики для средних и высших учебных заведений. 
Однако на сегодняшний день благодаря исследованиям существует способ повышения его КПД. Значительный рост производительности 
обеспечивается за счёт конструктивных доработок и использования материалов нового поколения. Также данное устройство позволяет 
освоить гидроэнергетические ресурсы каналов, малых рек и сбросных водоводов. Здесь использование ковшовых, поворотно-лопастных, 
радиально-осевых и диагональных турбин нецелесообразно из-за высоких затрат и низкой результативности.

Целью данного эссе является комплексный анализ исследований и актуальных применений сегнерова колеса. 
Предстоит провести оценку его энергетических характеристик по сравнению с подобными гидротурбинами. Важным направлением
является определение перспективных способов модернизации данного устройства. В ходе работы необходимо выполнить следующие
задачи:

\begin{enumerate}
    \item Изучить и представить историческое развитие сегнерова колеса и развитие его конструктивных особенностей.
    \item Проанализировать принцип работы и получить основные расчётные зависимости мощности, вращающего момента и реактивной силы.
    \item Сравнить эффективность различных модификаций сегнерова колеса.
    \item Оценить перспективность турбин данного типа.
\end{enumerate}

При подготовке были исследованы конструктивные особенности и энергетические характеристики. Проиллюстрированы различия в эксплуатации по сравнению с осевыми турбинами. В работе использованы научные публикации из проверенных баз данных, методы сравнительной оценки технических решений.